\documentclass[11pt]{article}

\usepackage[margin=1in]{geometry}
\usepackage{amsmath,amssymb,amsfonts}
\usepackage{booktabs}
\usepackage{multirow}
\usepackage{graphicx}
\usepackage{hyperref}
\usepackage{xcolor}
\hypersetup{
    colorlinks=true,
    linkcolor=black,
    citecolor=black,
    urlcolor=blue!50!black,
}
\usepackage{microtype}
\usepackage{authblk}

\title{Empirical Validation of the GROUND Ranking Framework on \\
       SP100 Weekly Credit Spreads (2020--2026)}

\author{Peter J.~Mercurio, PhD}
\affil{\small \texttt{peter.mercurio@alumni.utoronto.ca}}
\date{May 2026}


\begin{document}
\maketitle

\begin{abstract}
We present a six-year backtest of a weekly credit-spread strategy on the SP100
universe in which the GROUND ranking framework of \citet{mercurio2020}
selects, at each weekly entry, the top-$N$ candidate spreads from a Greek-filtered
candidate pool. We adopt a parsimonious intrinsic form of the GROUND ratio
($\Gamma = G / 3^{\,k\,D_{\mathrm{KL}}}$) and freeze the amplification factor
$k=20$ on a 2020--2024 in-sample window before scoring any 2025--2026 data.
On the in-sample window the strategy posts Sharpe 1.79 and a bankroll-independent
wager yield of 9.9\% at zero slippage; on the strict 2025--2026 holdout (350
trades, 70 weeks, never used in $k$-selection) it posts Sharpe 2.47 and yield
14.4\%. SPY buy-and-hold over the same windows posts Sharpe 0.79--1.11. Yield
degrades to 4.5--5.4\% on the full sample and 8.7\% on the holdout under a
conservative 3\textcent/leg post-hoc execution haircut. Single-factor baselines
that rank by Kelly-growth $G$ alone or KL-divergence $D_{\mathrm{KL}}$ alone are
reported for comparison. The data give direct empirical evidence that GROUND's
joint growth-and-divergence penalty produces a useful trade-selection signal in
a real options market with realistic execution costs.
\end{abstract}


\section{Introduction}

The GROUND ratio was introduced in \citet{mercurio2020} as a portfolio
ranking criterion that simultaneously rewards expected log-growth and
penalizes deviation from a uniform reference distribution. Concretely,
for a candidate trade with three discrete outcomes (full-win, partial,
full-loss) under empirical probabilities $(p, r_o, q)$ and an alpha
parameter $\alpha$ giving the partial-loss multiple, GROUND combines
\begin{itemize}
    \item the Kelly-optimal log-growth $G(p, q, r_o; \alpha)$, computed
        at the closed-form $w^\star$ solving the standard Kelly-with-partials
        first-order condition;
    \item the Kullback--Leibler divergence
        $D_{\mathrm{KL}}(\mathbf{p} \| \mathbf{u})$ from the uniform
        three-outcome distribution.
\end{itemize}
The intuition is straightforward: $G$ measures the trade's growth
\emph{per dollar at risk}, while $D_{\mathrm{KL}}$ measures the
\emph{informational} departure from a maximally-uncertain reference.
Trades with high $G$ but very high $D_{\mathrm{KL}}$ are
``sharp'' bets on a narrow outcome distribution and should be discounted;
trades with high $G$ and modest $D_{\mathrm{KL}}$ are diversified-edge
bets and should be preferred.

The original paper presented GROUND on synthetic three-outcome problems
and a closed-form Kelly setting. The work reported here is, to our
knowledge, the first end-to-end backtest of GROUND-driven trade selection
on a real options market, including realistic candidate generation,
filtering, and execution costs. We also adopt a simplified
\emph{intrinsic} form of GROUND that drops the per-week
minimum-divergence reference $R_b$ originally used to normalise the
denominator (Section~\ref{sec:ground-v3}). The intrinsic form is
algebraically simpler, computationally trivial, and empirically
indistinguishable from the reference-normalised form across our sweep.

The paper is organised as follows. Section~\ref{sec:strategy} describes
the data, candidate generation, filters, sizing, and slippage treatment.
Section~\ref{sec:ground-v3} states the GROUND v3 formula precisely and
notes our log-base canon. Section~\ref{sec:results} reports the in-sample
2020--2024 window, the strict 2025--2026 holdout, the full extended
2020--2026 sample, an SPY buy-and-hold benchmark, the $k$-sweep table, and
single-factor baselines that strip GROUND down to its $G$ and
$D_{\mathrm{KL}}$ components. Section~\ref{sec:discussion} discusses
limitations.


\section{Related Work}\label{sec:related}

This paper sits at the intersection of four literatures: the Kelly criterion
applied to derivatives, information-theoretic portfolio selection, systematic
short-volatility option strategies, and ranking-based trade selection. We
review each briefly to position the GROUND ratio's contribution.

\subsection{Kelly criterion in derivatives}

The Kelly criterion \citep{kelly1956} prescribes the bet-sizing fraction that
maximises expected log-wealth, and is the growth-optimal strategy in repeated
play. Its application to financial markets traces to \citet{thorp2006} and is
developed comprehensively in \citet{maclean2011}, including extensions to
fractional-Kelly sizing under estimation error and to multi-outcome bets such
as those characteristic of options spreads. \citet{vince1990} extended Kelly
to bounded-loss instruments where the gambler-style $f^\star$ does not apply
directly; credit spreads, with their bounded per-spread maximum loss, fall
naturally into this regime.

The GROUND framework extends this lineage in two ways. First, the closed-form
$w^\star$ used in the Kelly-growth term $G$ solves the first-order condition
for a three-outcome bet (full-win, partial, full-loss), which matches the
realised P\&L distribution of a held-to-expiry credit spread more faithfully
than a binary win/loss formulation. Second, single-factor ranking by $G$
alone is empirically unstable when probability estimates are noisy; the
GROUND ratio addresses this by penalising the divergence of the candidate's
outcome distribution from uniform, effectively discounting candidates whose
growth advantage relies on a sharp probability concentration. We report
$G$-only and $D_{\mathrm{KL}}$-only baselines in
Section~\ref{sec:baselines} as direct evidence.

\subsection{Information-theoretic portfolio selection}

The use of information-theoretic quantities in portfolio choice begins with
\citet{cover1991universal}, whose universal portfolios attain the growth rate
of the best constant-rebalanced portfolio in hindsight without distributional
assumptions, and \citet{algoet1988asymptotic}, who established the asymptotic
optimality of the log-optimal portfolio under general stationary processes.
\citet{cover2006} provide the canonical reference for the
Kullback--Leibler divergence and its role in characterising the cost of
distributional misspecification. The GROUND ratio's $D_{\mathrm{KL}}$ term
sits in this tradition: it measures the informational distance between the
candidate's empirical outcome distribution and a maximally-uncertain uniform
reference, and uses that distance as a regulariser on the
growth-maximisation objective.

The pair-trading antecedent \citep{mercurio2020} that introduced the GROUND
formulation used the divergence term in a comparative
(candidate-against-reference) form. The intrinsic v3 form adopted in the
present work removes the cross-candidate coupling and exposes
$D_{\mathrm{KL}}$ as a per-candidate regulariser, which is the more natural
formulation for streaming weekly selection.

\subsection{Systematic short-volatility strategies}

The realised behaviour of the strategy described here\,---\,collecting weekly
credit on bounded-loss vertical spreads, gated by a trend filter\,---\,is
adjacent to the systematic short-volatility literature. The CBOE PUT and
WPUT indices \citep{cboe_put} provide the canonical benchmark for short
cash-secured put strategies; \citet{israelovnielsen2014} document the
volatility risk premium that makes these strategies profitable in
expectation, and \citet{israelov2017} characterise their drawdown profile in
adverse regimes. Vertical credit spreads differ from short cash-secured puts
in that maximum loss is bounded per contract, which materially changes the
drawdown profile but preserves the underlying short-vol payoff structure.

Where the present work departs from this literature is in the
\emph{selection} step. Index strategies (PUT, WPUT) are mechanical:
short-the-ATM-put on a fixed schedule. The GROUND-driven strategy is
selective: at each weekly entry it ranks a Greek-filtered candidate pool of
up to several hundred spreads and chooses the top five. The empirical
question this paper addresses is whether the ranking step adds value over
mechanical or single-factor alternatives; Section~\ref{sec:results}
(in particular the single-factor baselines) provides the affirmative answer.

\subsection{Ranking-based trade selection}

Ranking-based selection is widely used in equity factor strategies
\citep{fama1992,asness2013} but is less well developed in derivatives.
\citet{goyenko2022} document that option mispricing signals can be ranked to
construct outperforming portfolios on the equity-options cross-section.
\citet{cao2023} apply machine-learning ranking to option returns directly.
The GROUND ratio differs from these in that it is a closed-form,
parameter-light ranker (one tuneable amplification factor $k$) derived from
first principles rather than fit to data; the only data-tuned parameter is
$k$, which we sweep on the in-sample window and freeze before the holdout
(Section~\ref{sec:k-sweep}).


\section{Strategy and Methodology}\label{sec:strategy}

\subsection{Universe and instrument}

The traded universe is the SP100 constituent set as of 2020-01-01, held
fixed throughout the backtest. We do not reconstruct point-in-time
membership; this introduces survivorship bias, since names that joined
the index after 2020-01-01 are excluded throughout, and names that
were dropped from the index during the window remain in the candidate
pool for their full backtest lifetime. The bias direction is ambiguous
but non-zero, and we make no attempt to quantify it.

At each weekly entry date we consider \emph{both} a vertical bull-put
spread and a vertical bear-call spread on each underlying that has a
Greek-feasible candidate pair, and let GROUND select between them.

Concretely, for each underlying and entry date the candidate generator
emits all $(K_{\text{short}}, K_{\text{long}})$ pairs satisfying:
\begin{enumerate}
    \item Days-to-expiry $\in [3, 8]$;
    \item Short-leg $|\Delta| \in [0.35, 0.65]$, target $|\Delta| = 0.50$;
    \item Per-share max-loss cap $\le \$5$;
    \item Net-credit $> 0$ (no net-debit fills);
    \item Credit-to-max-loss ratio $\ge 0.30$ (must collect at least
          30\textcent\ per dollar at risk; rejects under-priced wide-strike
          spreads with implausibly small credit relative to width);
    \item Short-leg open interest $\ge 10$ contracts (liquidity gate to
          screen out illiquid strikes with potentially stale or
          non-tradable quotes).
\end{enumerate}

\paragraph{Candidate pool size.}
Each weekly entry produces a candidate pool that exceeds the top-$N=5$
cap in every week of the backtest --- 1{,}300 trades over 260 in-sample
weeks is exactly $5\times$ the week count, indicating the cap binds 100\%
of the time. Across the in-sample window 77 distinct underlyings were
filled at least once. We do not log per-week pool sizes, but the
candidate pool is large enough every week that GROUND is doing real
selection, not picking from a pool of one or two.

\subsection{Macro filter}

A single macroeconomic filter is applied: SPY's close must be above its
trailing 100-day SMA for bull-put fills, and below for bear-call fills.
This is a binary regime gate, not a continuous tilt. Per-ticker regime
gating, gap filters, low-VIX filters, and theta/credit ratio filters
were all evaluated and found to be either redundant or
return-degrading; only the SPY regime gate is kept in the canonical
configuration.

\subsection{GROUND v3: intrinsic form}\label{sec:ground-v3}

The original GROUND formulation \citep{mercurio2020} ranks candidate
$a$ against a per-week minimum-divergence reference $R_b$:
\begin{equation}
    \Gamma_{a \mid b} \;=\; \frac{G_a - G_b}{1 + D_{\mathrm{KL}}^{(a)} - D_{\mathrm{KL}}^{(b)}}.
    \label{eq:ground-paper}
\end{equation}
This denominator regularisation is necessary in the original setting
(handling ties in minimum DKL) but introduces a cross-candidate
coupling that is awkward in a streaming weekly-selection setting and that
does not survive the Greek-driven candidate-pool variants we tested.

We instead adopt the \emph{intrinsic} v3 form:
\begin{equation}
    \Gamma_a \;=\; \frac{G_a}{3^{\,k\,D_{\mathrm{KL}}^{(a)}}},
    \qquad k = 20.
    \label{eq:ground-v3}
\end{equation}
The exponential denominator penalises high-divergence candidates
multiplicatively rather than additively. The amplification factor
$k$ controls how aggressively concentration risk is discounted; we
selected $k$ by sweeping $\{10, 15, 20, 25, 30\}$ on the 2020--2024
in-sample window. The full sweep is reported in
Table~\ref{tab:k-sweep}; $k=20$ sits on a Sharpe plateau of 1.73--1.79
spanning $k \in [20, 30]$. We chose $k=20$ from the low end of the
plateau and froze it before scoring any 2025+ data.

\paragraph{Logarithm base.} A credit spread has three discrete
outcomes (full-win, partial, full-loss). GROUND is therefore naturally
expressed in log-base-3 (trits), under which the uniform-outcome
entropy is $H(\mathbf{u}) = \log_3 3 = 1$ and divergence from uniform
satisfies $D_{\mathrm{KL}} \in [0, 1]$. The upper bound has direct
interpretation: $D_{\mathrm{KL}} = 1$ corresponds to a candidate
concentrating all probability on a single outcome. We compute $G$,
$D_{\mathrm{KL}}$, and the GROUND denominator $3^{k\,D_{\mathrm{KL}}}$
all in base 3, which keeps the exponent dimensionally consistent. Any
other log base produces an algebraically equivalent ranker at a
re-scaled $k$, so the choice of base is purely a unit convention.

\subsection{Sizing}

Two flat-sizing regimes are reported, $q \in \{1, 2\}$ contracts per
trade, with starting bankroll \$10{,}000. Under flat sizing the GROUND
ranking is bankroll-independent: the top-$N$ ordering at each weekly
entry depends only on per-spread economics, not on the running bankroll.
Consequently the $q=1$ and $q=2$ trade lists are identical; their
bankroll trajectories differ exactly by a factor of two (modulo the
floor at $\$0.01$ which is never hit in our runs).

\subsection{Slippage treatment}

We treat slippage as a \emph{post-hoc} P\&L haircut, not a
trade-selection input. Concretely, we run the backtest once at zero
slippage, recording the canonical trade list, and then apply a per-leg
haircut $s \in \{0\textcent, 1\textcent, 2\textcent, 3\textcent\}$ to
each fill. The haircut enters both the realised P\&L and the
``wagered'' figure used in the yield calculation (since slippage
inflates the effective max-loss per share).

This separation is important. If slippage is fed into trade selection,
the resulting trade list is a \emph{different} strategy at each
slippage level, and aggregate-statistics comparisons across slippage
become apples-to-oranges. With post-hoc treatment, we report a single
trade list under four execution-cost regimes.

\subsection{Bankroll-independent yield}

Conventional total-ROI is bankroll-dependent: a strategy that
turns \$10k into \$30k can be reported as ``200\% ROI'' but those
percentages do not translate to a different starting bankroll because
flat sizing does not compound. We instead report a sportsbook-style
\emph{yield}:
\begin{equation}
    \text{yield} \;=\; \frac{\sum_t q_t \cdot \text{pnl}_t \cdot 100}
                          {\sum_t q_t \cdot (\text{max\_loss}_t + 2s) \cdot 100} \;\times\; 100\%,
\end{equation}
which is invariant to starting bankroll and depends only on per-trade
economics. We retain conventional final-bankroll, Sharpe, and max
drawdown statistics for cross-comparability with other backtests.


\section{Results}\label{sec:results}

We report three windows. The \emph{in-sample} 2020--2024 window is the
window on which $k=20$ was selected from a sweep over $\{10, 15, 20, 25, 30\}$.
The \emph{holdout} 2025--2026 window contains 350 trades over 70 weeks
that were never used in $k$ selection or any other tuning; this is the
clean out-of-sample test. The \emph{extended sample} 2020--2026 window
combines the two; it inflates trade count and is reported only to make
clear what changes when 2025--2026 is appended to the train window
(and to anchor against the multi-slippage equity curves shown in the
companion HTML reports). The headline scientific claim rests on the
holdout, not the extended sample.

All numbers are at the canonical configuration described in
Section~\ref{sec:strategy}, with GROUND threshold zero (every selected
trade is taken), top-$N=5$ per week, and base-3 canon. Reported
yields are gross of commissions; a commission-inclusive row is added
in Section~\ref{sec:results-commissions}.

\begin{table}[h]
\centering
\small
\begin{tabular}{lrrrrrr}
\toprule
window           & Trades & Final (0\textcent) & Sharpe & Max DD & Yield (0\textcent) & Yield (3\textcent) \\
\midrule
\multicolumn{7}{l}{\emph{$q=1$ contract per trade}} \\
in-sample 2020--2024   & 1{,}300 & \$24{,}973 & 1.79 & $-9.3\%$  & 9.92\%  & 4.52\% \\
holdout 2025--2026     & 350     & \$15{,}755 & 2.47 & $-7.4\%$  & 14.37\% & 8.67\% \\
extended 2020--2026    & 1{,}650 & \$30{,}729 & 1.94 & $-9.3\%$  & 10.85\% & 5.39\% \\
\midrule
\multicolumn{7}{l}{\emph{$q=2$ contracts per trade --- selection-identical to $q=1$, P\&L doubled}} \\
in-sample 2020--2024   & 1{,}300 & \$39{,}946 & 1.79 & $-12.2\%$ & 9.92\%  & 4.52\% \\
extended 2020--2026    & 1{,}650 & \$51{,}457 & 1.94 & $-12.2\%$ & 10.85\% & 5.39\% \\
\bottomrule
\end{tabular}
\caption{Headline statistics. The holdout 2025--2026 row is the strict
OOS evaluation: $k=20$ was frozen from the in-sample window before any
2025+ data was scored. The extended-sample row appends 2025--2026 to
the in-sample window without re-tuning. Yield is bankroll-independent
realised-P\&L over wagered. The 0\textcent\ column is the
selection-canonical (mid-mid) baseline; the 3\textcent\ column applies
a per-leg post-hoc haircut as a conservative execution-cost regime.}
\label{tab:variants}
\end{table}

\subsection{Holdout vs.\ in-sample}

On the strict 2025--2026 holdout, Sharpe is 2.47 and yield is 14.4\%
at zero slippage. Both are higher than the in-sample window, not lower.
We do not over-interpret this: the holdout is short (70 weeks) and
2025--2026 was a generally favourable regime for the underlying SPY
trend filter (SPY closed above its 100-day SMA most of the holdout).
The fact that the strategy does not \emph{collapse} on unseen data is
the load-bearing claim; the apparent improvement is consistent with
sampling noise plus regime tailwind.

\subsubsection{Confidence intervals on Sharpe}\label{sec:sharpe-ci}

Sharpe ratio estimates from short samples have wide confidence intervals,
and we report them here to forestall over-reading of the headline
numbers. \citet{lo2002} derives the asymptotic distribution of the
estimated Sharpe ratio under the assumption of i.i.d.\ Gaussian returns:
\begin{equation}
\widehat{\mathrm{SR}}_w \overset{a}{\sim}
\mathcal{N}\!\left(\mathrm{SR}_w,\;
\frac{1 + \tfrac{1}{2}\mathrm{SR}_w^2}{T}\right),
\label{eq:lo-sharpe}
\end{equation}
where $T$ is the number of return observations. Credit-spread weekly
P\&L, however, is mechanically bounded above by the credit collected
and below by the max-loss, which produces a heavily left-skewed return
distribution that violates the Gaussian assumption underlying
\eqref{eq:lo-sharpe}. \citet{mertens2002} provides the
higher-moment-adjusted variance that is valid under i.i.d.\ but
non-Gaussian returns (also reproduced in \citet{opdyke2007}):
\begin{equation}
\mathrm{Var}(\widehat{\mathrm{SR}}_w)
  = \frac{1 + \tfrac{1}{2}\mathrm{SR}_w^2 \;-\; \gamma_3 \mathrm{SR}_w
            \;+\; \tfrac{1}{4}\gamma_4 \mathrm{SR}_w^2}{T},
\label{eq:mertens-sharpe}
\end{equation}
where $\gamma_3$ is the sample skewness and $\gamma_4$ the sample
excess kurtosis of the weekly returns. For the canonical
short-volatility credit-spread strategy described here, $\gamma_3 < 0$
(frequent small wins, occasional large losses) and $\gamma_4 > 0$ (fat
tails). Both terms inflate $\mathrm{Var}(\widehat{\mathrm{SR}})$:
$-\gamma_3 \mathrm{SR}$ is positive because $\gamma_3$ is negative and
the Sharpe is positive, and $\gamma_4 \mathrm{SR}^2 / 4$ is positive
because $\gamma_4$ is positive. Mertens (2002) therefore produces
\emph{wider} confidence intervals than Lo (2002) on this strategy,
which is the conservative direction.

Annualised standard errors are obtained by scaling
\eqref{eq:mertens-sharpe} by $\sqrt{52}$, and 95\% intervals by
$\widehat{\mathrm{SR}}_{\mathrm{ann}} \pm 1.96 \cdot
\mathrm{SE}(\widehat{\mathrm{SR}}_{\mathrm{ann}})$. Both Lo and Mertens
intervals are reported in Table~\ref{tab:sharpe-ci}.

\begin{table}[h]
\centering
\footnotesize
\begin{tabular}{lrrrrrrrr}
\toprule
window & $T$ & SR & $\gamma_3$ & $\gamma_4$ & Lo SE & Lo 95\% CI & Mertens SE & Mertens 95\% CI \\
\midrule
in-sample 2020--2024 & 260 & 1.79 & $-0.81$ & $+0.64$ & 0.45 & $[0.90,\,2.68]$ & 0.50 & $[0.81,\,2.76]$ \\
holdout 2025--2026   &  70 & 2.47 & $-0.81$ & $+0.69$ & 0.89 & $[0.73,\,4.21]$ & 1.00 & $[0.50,\,4.44]$ \\
extended 2020--2026  & 330 & 1.94 & $-0.80$ & $+0.62$ & 0.40 & $[1.15,\,2.73]$ & 0.45 & $[1.06,\,2.81]$ \\
\midrule
SPY 2020--2024       & 261 & 0.79 & $-0.49$ & $+5.56$ & 0.45 & $[-0.09,\,1.66]$ & 0.46 & $[-0.12,\,1.69]$ \\
SPY 2025--2026       &  72 & 1.11 & $-0.55$ & $+3.36$ & 0.85 & $[-0.57,\,2.78]$ & 0.90 & $[-0.65,\,2.87]$ \\
SPY 2020--2026       & 333 & 0.84 & $-0.50$ & $+5.47$ & 0.40 & $[0.07,\,1.62]$ & 0.41 & $[0.04,\,1.65]$ \\
\bottomrule
\end{tabular}
\caption{Annualised Sharpe ratios with iid-Gaussian Lo (2002) standard
errors and iid-non-Gaussian Mertens (2002) standard errors that
incorporate sample skewness $\gamma_3$ and excess kurtosis $\gamma_4$.
Strategy returns are weekly P\&L scaled by the fixed \$10{,}000
starting bankroll. SPY benchmark returns are weekly (W-Friday) total
returns derived from Yahoo Finance dividend-adjusted closes. The
strategy's negative skew (–0.80) widens its CIs by roughly 11\% under
Mertens; SPY's fatter tails (excess kurtosis 5.5) widen its CIs by
only 2--6\% because the SPY Sharpe is small enough that the
$\gamma_4 \mathrm{SR}^2/4$ term is correspondingly small.}
\label{tab:sharpe-ci}
\end{table}

The Mertens intervals are the appropriate primary reference under the
return distribution this strategy actually produces. Three observations
follow.

First, the holdout's Mertens interval $[0.50,\,4.44]$ is consistent
with strategies ranging from ``moderately better than SPY'' to
``implausibly good.'' We report the 2.47 point estimate as the
headline but note that 70 weekly observations is insufficient to
discriminate sharply between these hypotheses. The in-sample interval
$[0.81,\,2.76]$ is tighter and materially exceeds SPY's matched-window
interval $[-0.12,\,1.69]$, but since $k$ was selected on this window
the in-sample Sharpe is not a clean estimate of the strategy's true
Sharpe.

Second, the extended-sample row is the most defensible single number
for externally communicating the strategy's risk-adjusted return: it
has the largest $T$, the tightest confidence interval, and combines
the $k$-selected window with the holdout. Its Mertens 95\% lower
bound of 1.06 exceeds SPY's extended-sample point estimate of 0.84,
though SPY's upper bound 1.65 overlaps the strategy's lower bound; the
windows are not independent, so this comparison is descriptive rather
than a formal test of difference.

Third, both Lo and Mertens assume independence of weekly observations.
Weekly aggregation removes within-week selection correlation but does
not eliminate autocorrelation across weeks under regime shifts. We do
not adjust for serial dependence here; the practical implication is
that all reported intervals should still be read as lower bounds on
uncertainty.

\subsection{SPY buy-and-hold benchmark}

\begin{table}[h]
\centering
\small
\begin{tabular}{lrrr}
\toprule
window & Sharpe & Total ROI & Max DD \\
\midrule
SPY 2020--2024 & 0.79 & 99.1\%  & $-31.8\%$ \\
SPY 2025--2026 & 1.11 & 26.6\%  & $-16.9\%$ \\
SPY 2020--2026 & 0.84 & 150.8\% & $-31.8\%$ \\
\bottomrule
\end{tabular}
\caption{SPY buy-and-hold total return over the matching evaluation
windows, derived from Yahoo Finance dividend-adjusted closes resampled
to W-Friday weekly returns. Sharpe is annualised at $\sqrt{52}$.}
\label{tab:spy-benchmark}
\end{table}

The strategy's Sharpe exceeds SPY's on every window, with a particularly
large margin on the in-sample window (1.79 vs.\ 0.79). On the holdout,
the margin tightens to 2.47 vs.\ 1.11 because SPY's holdout performance
is itself unusually strong. Drawdown is $-7$ to $-12$\% for the strategy
versus $-19$ to $-34$\% for SPY across windows, reflecting the
short-DTE credit-spread structure: max-loss is bounded per spread.

\subsection{Yield robustness to slippage}

On the in-sample window, yield falls from 9.92\% (0\textcent) to 4.52\%
(3\textcent) --- a 54\% haircut. On the holdout it falls from 14.37\% to
8.67\% --- a 40\% haircut. The 3\textcent/leg figure is intentionally
conservative; on liquid SP100 weeklies, 1--2\textcent/leg better
characterises realised execution. At those midpoint costs, in-sample
yield holds 8.13\% (1\textcent) and 6.32\% (2\textcent).

\subsection{Commission-inclusive yield}\label{sec:results-commissions}

The numbers above are gross of commissions. At a typical retail rate of
\$0.65/contract/leg and a credit-spread structure where positions are
held to expiry (entry only, no closing fill), commissions add roughly
\$1.30/contract. On the in-sample 1{,}300-trade $q=1$ window that is
\$1{,}690 of commission, which reduces 0\textcent\ yield from 9.92\%
to 8.80\%. The shift is approximately 1\textcent/leg in
slippage-equivalent terms, so a reader who wants commission-inclusive
yields can read the 1\textcent\ slippage column with a small additional
discount: roughly 8.1\% in-sample, 12.4\% holdout.

\subsection{$k$-sweep on the in-sample window}\label{sec:k-sweep}

Table~\ref{tab:k-sweep} reports the ranking-amplification sweep
$k \in \{10, 15, 20, 25, 30\}$ on the in-sample window, all at $q=1$.
This was the data on which $k=20$ was selected.

\begin{table}[h]
\centering
\small
\begin{tabular}{rrrrrr}
\toprule
$k$ & Trades & Final & Sharpe & Max DD & Yield \\
\midrule
10 & 1{,}300 & \$21{,}761 & 1.35 & $-7.5\%$  & 7.71\% \\
15 & 1{,}300 & \$23{,}324 & 1.55 & $-10.3\%$ & 8.72\% \\
\textbf{20} & \textbf{1{,}300} & \textbf{\$24{,}973} & \textbf{1.79} & \textbf{$-9.3\%$} & \textbf{9.92\%} \\
25 & 1{,}300 & \$24{,}766 & 1.73 & $-8.9\%$  & 9.93\% \\
30 & 1{,}300 & \$25{,}853 & 1.79 & $-8.7\%$  & 10.77\% \\
\bottomrule
\end{tabular}
\caption{$k$-sweep on the in-sample 2020--2024 window, $q=1$, base-3 canon.
The Sharpe ratio rises from 1.35 at $k=10$ to a 1.73--1.79 plateau over
$k \in [20, 30]$. $k=20$ was chosen from the low end of the plateau and
frozen before any 2025+ data was scored. Yield continues to drift
upward at $k=30$ but Sharpe does not, suggesting marginal gains from
larger $k$ come from a few high-yield, higher-variance trades.}
\label{tab:k-sweep}
\end{table}

\subsection{Single-factor baselines}\label{sec:baselines}

To test whether GROUND adds value over its components, we re-rank with
two stripped-down rankers: (i) $G$-only (rank by Kelly-growth alone,
ignoring $D_{\mathrm{KL}}$), and (ii) $D_{\mathrm{KL}}$-only (rank by
lowest divergence, ignoring $G$). Both run on the in-sample 2020--2024
window, $q=1$, with all other configuration unchanged.

\begin{table}[h]
\centering
\small
\begin{tabular}{lrrrrr}
\toprule
ranker & Trades & Final & Sharpe & Max DD & Yield \\
\midrule
\textbf{GROUND ($k=20$)}   & 1{,}300 & \$24{,}973 & \textbf{1.79} & $-9.3\%$  & \textbf{9.92\%} \\
$G$ alone (Kelly-growth)   & 1{,}300 & \$20{,}555 & 1.22          & $-9.0\%$  & 6.92\% \\
$D_{\mathrm{KL}}$ alone    & 1{,}300 & \$18{,}704 & 0.83          & $-24.5\%$ & 6.86\% \\
\bottomrule
\end{tabular}
\caption{Single-factor baselines, in-sample 2020--2024, $q=1$.
$G$-only ranks by Kelly-growth and ignores $D_{\mathrm{KL}}$;
$D_{\mathrm{KL}}$-only ranks by lowest divergence and ignores $G$.
Both underperform the joint GROUND ranker on Sharpe and yield. The
$D_{\mathrm{KL}}$-only ranker is particularly punishing on drawdown
($-24.5\%$ vs.\ $-9.3\%$ for GROUND), consistent with the intuition
that selecting the \emph{flattest} probability distribution does not
control growth.}
\label{tab:baselines}
\end{table}

\subsection{Sizing scales linearly}

The trade list is identical between $q=1$ and $q=2$ at the same date
window; final bankroll scales linearly within the floor-free regime,
while yield, Sharpe, and drawdown are sizing-invariant by construction.


\section{Discussion}\label{sec:discussion}

\subsection{Why GROUND vs.\ simpler rankers}

We compare GROUND against two single-factor baselines (Section~\ref{sec:baselines}):
$G$-only, which ranks by Kelly-growth and ignores $D_{\mathrm{KL}}$,
and $D_{\mathrm{KL}}$-only, which ranks by lowest divergence and ignores
$G$. The full GROUND ranker outperforms both on Sharpe and yield on the
in-sample window (Table~\ref{tab:baselines}). The two components are
individually weak signals; their multiplicative combination via
Equation~\ref{eq:ground-v3} is where the empirical edge appears.

\subsection{Limitations}

\begin{enumerate}
    \item \emph{Survivorship bias.} The SP100 constituent list is held
        fixed at the 2020-01-01 composition. Names that joined the
        index after that date are excluded, and names that left the
        index during the window remain in the candidate pool. The
        bias direction is ambiguous and unquantified.
    \item \emph{Holdout length.} The strict 2025--2026 holdout is
        70 weeks, which is enough to test for collapse but not enough
        to characterise the distribution of returns. Sharpe estimates
        from a 70-week sample have wide confidence intervals.
    \item \emph{Greek-only candidate generator.} We do not blend
        implied volatility with realised volatility or use
        drift-adjusted probabilities in the canonical configuration.
        Both were tested as toggles in earlier iterations and found to
        be non-improving on Sharpe under the GROUND ranker.
    \item \emph{Fill assumption.} Mid-mid quotes are taken as
        executable, with the post-hoc per-leg haircut serving as the
        execution-cost regime.
    \item \emph{Commissions.} Section~\ref{sec:results-commissions}
        reports a commission-inclusive estimate; for a retail account
        at $\sim\$0.65$/contract per leg this shifts yields by roughly
        one slippage column.
    \item \emph{Single macro filter.} The SPY-100d-SMA gate is binary
        and global. A more sophisticated regime model would likely
        improve performance but would also increase overfitting risk
        and was deliberately excluded from the canonical configuration.
\end{enumerate}

\subsection{Conclusion}

Under the intrinsic v3 form $\Gamma = G / 3^{\,k\,D_{\mathrm{KL}}}$
with $k=20$ in trits (log-base-3), the GROUND ranker yields Sharpe 1.79 on the
2020--2024 in-sample window and Sharpe 2.47 on the 350-trade,
70-week strict 2025--2026 holdout, against SPY buy-and-hold Sharpes
of 0.79 and 1.11 over the same windows. Under a 3\textcent/leg
post-hoc execution haircut, yield falls from 9.9\% to 4.5\% in-sample
and from 14.4\% to 8.7\% on the holdout. Single-factor baselines that
strip out either $G$ or $D_{\mathrm{KL}}$ underperform the combined
ranker (Section~\ref{sec:baselines}).


\appendix

\section{Closed-form Kelly-growth score with a partial outcome}
\label{app:kelly-foc}

Section~\ref{sec:strategy} references the Kelly-growth score $G$ used as
the numerator of the GROUND ratio, computed at a closed-form
$w^\star$. We derive that fraction here.

\paragraph{Score, not sizing.}
The Kelly-growth $G$ used in this work is a \emph{ranking score}, not a
prescription for how many contracts to trade. The actual sizing rule is
flat ($q=1$ or $q=2$ contracts per spread, capped at
\texttt{MAX\_CONTRACTS}); $w^\star$ never enters portfolio construction.
The score lives on a normalised three-outcome bet whose payoffs are
$\{+1, +\alpha, -1\}$, parameter-free apart from the partial multiplier
$\alpha$. This normalisation lets us compare credit spreads of different
sizes on a common scale: a thin-credit spread and a fat-credit spread
both map to the same outcome geometry and are differentiated only
through their probabilities $(p, r_o, q)$ and through $\alpha$.

\paragraph{Setup.}
Consider a wager with three outcomes:
\begin{itemize}
  \item full win, with probability $p$ and per-dollar payoff $+1$;
  \item partial outcome, with probability $r_o$ and per-dollar payoff
        $+\alpha$, where $\alpha \in [-1, 1]$ is the partial multiplier;
  \item full loss, with probability $q = 1 - p - r_o$ and per-dollar
        payoff $-1$.
\end{itemize}
The canonical configuration uses $\alpha = -0.5$, taken from
\citet{mercurio2020}, equation~28: the average partial-outcome return on
a held-to-expiry credit spread that expires between the two strikes,
expressed in normalised units. The probability $q$ is decomposed from
the short-leg ITM probability via $q = \Delta_{\text{short}} \cdot
\texttt{LOSS\_FACTOR}$ with $\texttt{LOSS\_FACTOR} = 0.60$ (canonical),
splitting the ITM mass between full-loss and partial outcomes.

\paragraph{Growth function.}
If a fraction $w \in [0, 1]$ of the (notional) bankroll is wagered, the
bankroll multiplier in each outcome is $1 + w$, $1 + w\alpha$, and
$1 - w$ respectively. The expected log-growth is
\begin{equation}
G(w; p, r_o, q, \alpha)
  = p \log(1 + w) + r_o \log(1 + w\alpha) + q \log(1 - w).
\label{eq:growth}
\end{equation}

Differentiating with respect to $w$ and setting the derivative to zero
gives the first-order condition
\begin{equation}
\frac{p}{1 + w}
+ \frac{r_o\, \alpha}{1 + w \alpha}
- \frac{q}{1 - w} = 0.
\label{eq:foc}
\end{equation}

\paragraph{Closed form.}
Multiplying both sides of \eqref{eq:foc} by $(1+w)(1+w\alpha)(1-w)$ and
collecting powers of $w$ yields the quadratic
\begin{equation}
A w^2 + B w + C = 0,
\label{eq:quadratic}
\end{equation}
with coefficients
\begin{align*}
A &= -\alpha\,(p + r_o + q) = -\alpha, \\
B &= \alpha (p - q) - (p + q), \\
C &= p + r_o\, \alpha - q.
\end{align*}
The first equality in $A$ uses $p + r_o + q = 1$. The growth-optimal
fraction $w^\star$ is the unique root of \eqref{eq:quadratic} in
$(0, 1)$:
\begin{equation}
w^\star = \frac{-B - \sqrt{B^2 - 4AC}}{2A}.
\label{eq:wstar}
\end{equation}
For the canonical $\alpha = -0.5$, $A > 0$ and the in-unit root is
uniquely $w^\star$ above. The clearing of the $(1+w\alpha)$ factor
introduces a spurious second root outside $(0, 1)$ which is discarded.
The growth score reported in $G$ is then
\begin{equation}
G = p \log_3(1 + w^\star) + r_o \log_3(1 + w^\star \alpha)
    + q \log_3(1 - w^\star),
\label{eq:G-used}
\end{equation}
computed in base 3 as discussed in Section~\ref{sec:ground-v3} so that
$G$ shares units with $D_{\mathrm{KL}}$ in the GROUND ratio.

\paragraph{Boundary cases.}
When $\alpha = -1$ the partial outcome is indistinguishable from a full
loss; combining $r_o$ into $q$ recovers the standard symmetric binary
Kelly $w^\star = p - (r_o + q) = 2p - 1$. When $\alpha = 1$ the partial
outcome is indistinguishable from a full win; the analogous combination
with $p$ gives $w^\star = (p + r_o) - q = 1 - 2q$. When $\alpha = 0$ the
partial outcome is a wash and \eqref{eq:foc} reduces to a binary Kelly
on the renormalised pair $(p, q)$ with $w^\star = (p - q) / (p + q)$.

\paragraph{Empirical probability estimation.}
The probabilities $(p, r_o, q)$ are estimated from a Black--Scholes
risk-neutral distribution under the canonical configuration, with
short-leg implied volatility used as the volatility input and
$\texttt{LOSS\_FACTOR}$ splitting the ITM mass. We tested
realised-volatility, drift-adjusted, and blended IV/RV estimators in
earlier iterations and found no Sharpe improvement under the GROUND
ranker; this robustness to the probability-estimation choice is
consistent with the intuition that GROUND's $D_{\mathrm{KL}}$ regulariser
absorbs much of the estimator's distributional error.


\bibliographystyle{plainnat}
\begin{thebibliography}{9}

\bibitem[Mercurio et~al.(2020)]{mercurio2020}
P.\,J.~Mercurio, Y.~Wu, and Q.~Xie.
\newblock An entropic approach for pair trading.
\newblock \emph{Entropy}, 22(8):805, 2020.
\newblock \doi{10.3390/e22080805}.

\bibitem[Kelly(1956)]{kelly1956}
J.\,L.~Kelly~Jr.
\newblock A new interpretation of information rate.
\newblock \emph{Bell System Technical Journal}, 35(4):917--926, 1956.

\bibitem[Cover and Thomas(2006)]{cover2006}
T.\,M.~Cover and J.\,A.~Thomas.
\newblock \emph{Elements of Information Theory}.
\newblock Wiley-Interscience, 2nd edition, 2006.

\bibitem[MacLean et~al.(2011)]{maclean2011}
L.\,C.~MacLean, E.\,O.~Thorp, and W.\,T.~Ziemba (eds.).
\newblock \emph{The Kelly Capital Growth Investment Criterion}.
\newblock World Scientific, 2011.

\bibitem[Thorp(2006)]{thorp2006}
E.\,O.~Thorp.
\newblock The Kelly criterion in blackjack, sports betting, and the stock market.
\newblock In \emph{Handbook of Asset and Liability Management}, volume~1,
pages 385--428. 2006.

\bibitem[Vince(1990)]{vince1990}
R.~Vince.
\newblock \emph{Portfolio Management Formulas: Mathematical Trading Methods
for the Futures, Options, and Stock Markets}.
\newblock Wiley, 1990.

\bibitem[Cover(1991)]{cover1991universal}
T.\,M.~Cover.
\newblock Universal portfolios.
\newblock \emph{Mathematical Finance}, 1(1):1--29, 1991.

\bibitem[Algoet and Cover(1988)]{algoet1988asymptotic}
P.\,H.~Algoet and T.\,M.~Cover.
\newblock Asymptotic optimality and asymptotic equipartition properties
of log-optimum investment.
\newblock \emph{Annals of Probability}, 16(2):876--898, 1988.

\bibitem[CBOE(2024)]{cboe_put}
Chicago Board Options Exchange.
\newblock CBOE S\&P 500 PutWrite Index (PUT) methodology.
\newblock \url{https://www.cboe.com/us/indices/dashboard/put/}, 2024.

\bibitem[Israelov and Nielsen(2014)]{israelovnielsen2014}
R.~Israelov and L.\,N.~Nielsen.
\newblock Covered call strategies: One fact and eight myths.
\newblock \emph{Financial Analysts Journal}, 70(6):23--31, 2014.

\bibitem[Israelov(2017)]{israelov2017}
R.~Israelov.
\newblock Pathetic protection: The elusive benefits of protective puts.
\newblock \emph{Journal of Alternative Investments}, 19(3):38--56, 2017.

\bibitem[Fama and French(1992)]{fama1992}
E.\,F.~Fama and K.\,R.~French.
\newblock The cross-section of expected stock returns.
\newblock \emph{Journal of Finance}, 47(2):427--465, 1992.

\bibitem[Asness et~al.(2013)]{asness2013}
C.\,S.~Asness, T.\,J.~Moskowitz, and L.\,H.~Pedersen.
\newblock Value and momentum everywhere.
\newblock \emph{Journal of Finance}, 68(3):929--985, 2013.

% TODO: verify exact citation for Goyenko/Zhang option-mispricing ranking work
\bibitem[Goyenko and Zhang(2022)]{goyenko2022}
R.~Goyenko and C.~Zhang.
\newblock The joint cross-section of option and stock returns predictability.
\newblock Working paper, 2022.

% TODO: verify exact citation for Cao et al. ML option-return ranking work
\bibitem[Cao et~al.(2023)]{cao2023}
J.~Cao, B.~Han, Q.~Tong, and X.~Zhan.
\newblock Option return predictability with machine learning.
\newblock Working paper, 2023.

\bibitem[Lo(2002)]{lo2002}
A.\,W.~Lo.
\newblock The statistics of Sharpe ratios.
\newblock \emph{Financial Analysts Journal}, 58(4):36--52, 2002.

\bibitem[Mertens(2002)]{mertens2002}
E.~Mertens.
\newblock Comments on variance of the IID estimator in Lo (2002).
\newblock Unpublished working paper.

\bibitem[Opdyke(2007)]{opdyke2007}
J.\,D.~Opdyke.
\newblock Comparing Sharpe ratios: So where are the $p$-values?
\newblock \emph{Journal of Asset Management}, 8(5):308--336, 2007.

\end{thebibliography}

\end{document}
